\section{Notation and technical preliminaries}
\label{app:prelim}

For an episodic MDP, $h\in[H]$ indexes the stage, $s\in\cS$ the state,
$a\in\cA$ the action, and $z=(s,a)\in\cZ$. Rewards lie in $[0,1]$.
Policies and learning algorithms may be randomized and history dependent. The
kernel $k$ is positive semidefinite and has RKHS $\cH$. Its maximum information
gain at regularization $\lambda>0$ is
\begin{equation}
 \gamma_T(k;\lambda)=
 \sup_{z_1,\ldots,z_T\in\cZ}
 \frac12\log\det(I_T+\lambda^{-1}K_T),
 \qquad [K_T]_{ij}=k(z_i,z_j).
 \label{eq:gamma-app}
\end{equation}
We allow repeated points in the supremum. This convention has no effect on the
arguments.

We use the following elementary pullback fact.

\begin{lemma}[Pullback RKHS]
\label{lem:pullback}
Let $q:\widetilde\cZ\to\cZ$ be surjective and let
$\widetilde k(\tilde z,\tilde z')=k(q(\tilde z),q(\tilde z'))$. Then
\begin{equation}
 \widetilde\cH=\{f\circ q:f\in\cH\},
 \qquad \|f\circ q\|_{\widetilde\cH}=\|f\|_{\cH}.
 \label{eq:pullback-isometry}
\end{equation}
Also $\gamma_T(\widetilde k;\lambda)=\gamma_T(k;\lambda)$.
\end{lemma}

\begin{proof}
If $\phi:\cZ\to\mathcal F$ is a feature map for $k$, then
$\widetilde\phi=\phi\circ q$ is a feature map for $\widetilde k$. Surjectivity
means the closed linear spans of $\phi(\cZ)$ and
$\widetilde\phi(\widetilde\cZ)$ coincide. Their canonical feature-space
representations therefore give equation~\eqref{eq:pullback-isometry}. Every
$T$-point Gram matrix in the refined domain is the Gram matrix of its projected
points. Conversely, every original sequence has a lift because $q$ is surjective.
The two suprema in equation~\eqref{eq:gamma-app} are equal.
\end{proof}

\section{Proofs for exact state refinement}
\label{app:refinement}

We first state the full MDP construction. For clarity, take a layered MDP, so a
state cannot occur at two different stages. This loses no generality because the
stage can be appended to the state. For every $s\in\cS_h$, choose a finite integer
$m_s\geq1$ and let
\begin{equation}
 \widetilde\cS_h=\{(s,i):s\in\cS_h,\ i\in[m_s]\}.
\end{equation}
If the original initial distribution is $\rho$, set
$\widetilde\rho(s,i)=\rho(s)/m_s$. Define
\begin{align}
 \widetilde r_h((s,i),a)&=r_h(s,a),
 \label{eq:refined-reward-app}\\
 \widetilde P_h((s',j)\mid(s,i),a)
 &=\frac{P_h(s'\mid s,a)}{m_{s'}}.
 \label{eq:refined-p-app}
\end{align}
The projection $q(s,i)=s$ is a probabilistic bisimulation map.

\subsection{Equality of decision and learning experiments}

Given an original policy $\pi$, define its lift by ignoring clone indices. Summing
equation~\eqref{eq:refined-p-app} over $j$ proves inductively that the projected
trajectory has the original law. Rewards agree pathwise under projection.

The converse must account for a refined policy that uses clone indices. An agent in
the original MDP can generate them privately. At the initial state $s_1$, draw
$I_1$ uniformly from $[m_{s_1}]$. After observing each subsequent original state
$s_h$, draw $I_h$ independently and uniformly from $[m_{s_h}]$. Feed the synthetic
history $((s_1,I_1),a_1,\ldots,(s_h,I_h))$ to the refined policy. Conditional on
the projected history, equation~\eqref{eq:refined-p-app} makes the actual clone
index in the refined MDP exactly independent and uniform. Induction therefore
couples the full reward and projected-state laws.

The same simulation applies across episodes and to a learning algorithm whose next
policy depends on all earlier observations. One direction projects observations.
The other supplies private clone indices. Thus the sets of attainable reward laws
are identical. Optimal values are identical as a special case. Since regret is a
function of the selected-policy values and optimal values, the minimax regret is
identical for every interaction budget. This proves part 1 of
Theorem~\ref{thm:refinement}.

\subsection{Coordinate norm and information gain}

At a transition into $s'$, each refined coordinate function is
\begin{equation}
 \widetilde p_{s',j}((s,i),a)
 =\frac{1}{m_{s'}}p_{s'}(s,a).
\end{equation}
Lemma~\ref{lem:pullback} and homogeneity of the RKHS norm give
\begin{equation}
 \|\widetilde p_{s',j}\|_{\widetilde\cH}
 =\frac{\|p_{s'}\|_{\cH}}{m_{s'}}.
\end{equation}
Taking the supremum proves equation~\eqref{eq:coordinate-shrink}.
The information-gain claim is Lemma~\ref{lem:pullback}. This completes the proof of
Theorem~\ref{thm:refinement}.

\subsection{Proof of Corollary~\ref{cor:vacuity}}

On a finite domain, strict positive definiteness means the Gram matrix is
invertible. Every scalar function $f$ belongs to the RKHS and has finite norm
$\|f\|_{\cH}^2=f^\top K^{-1}f$. There are finitely many stages and target states.
Choose
\begin{equation}
 m_y\geq\max_h\frac{\|P_h(y\mid\cdot)\|_{\cH}}{\epsilon}
\end{equation}
and round upward. Theorem~\ref{thm:refinement} gives the coordinate bound and all
invariances. A pulled-back reward has the same norm by
Lemma~\ref{lem:pullback}.

\section{Proof of the linear-regret theorem}
\label{app:lower}

This section gives one domain and one kernel shared by every interaction budget.
The additional geometric detail ensures that all transition coordinates, including
their values on state-action pairs that are unreachable at a given stage, are
restrictions of Mat\'ern RKHS functions.

\subsection{A fixed compact state-action domain}

Let
\begin{equation}
 \cS=\{s_0,g_\infty,b_\infty\}
 \cup\{g_n,b_n:n\in\N\},
 \qquad \cA=[0,1].
\end{equation}
Assign state coordinates
\begin{align}
 \eta(s_0)&=0,\\
 \eta(g_\infty)&=\tfrac13,
 &\eta(g_n)&=\tfrac13+\frac{1}{100(n+1)},\\
 \eta(b_\infty)&=\tfrac23,
 &\eta(b_n)&=\tfrac23+\frac{1}{100(n+1)}.
\end{align}
Embed $z=(s,a)$ as $\iota(z)=(a,\eta(s))\in[0,1]^2$. The image $\iota(\cZ)$
is compact because both limit-state lines are included. Let $\kappa_\nu$ be the
unit-variance Mat\'ern kernel of smoothness $\nu>1/2$ on $[0,1]^2$ and define
\begin{equation}
 k(z,z')=\kappa_\nu(\iota(z),\iota(z')).
 \label{eq:matern-pullback}
\end{equation}
The embedding is injective, so this is simply the restriction of a fixed Mat\'ern
kernel to a compact subdomain.

Choose $q\in C^\infty([0,1])$ that equals one on the set of good-state
coordinates and zero on the initial and bad-state coordinates. Such a function
exists because these compact sets are separated. The function
$Q(a,y)=q(y)$ has a smooth compactly supported extension to $\R^2$. Mat\'ern RKHSs
on bounded Lipschitz domains are norm equivalent to Sobolev spaces of finite order
\citep{wendland2005scattered,kanagawa2018review}. Hence the restriction of $Q$ to
$\iota(\cZ)$ has a finite RKHS norm $C_q$. Set
\begin{equation}
 \alpha=\min\{1,C_q^{-1}\}>0,
 \qquad r_1\equiv0,
 \qquad r_2(s,a)=\alpha q(\eta(s)).
 \label{eq:fixed-reward}
\end{equation}
The reward is fixed, known, in $[0,1]$, and has RKHS norm at most one. This reward
smoothness is not needed by the open problem, but shows the lower bound survives
the reward part of the stronger coordinate assumptions used elsewhere.

\subsection{The transition family}

Fix $T\geq2$, set $M=2T$, and choose pairwise disjoint open intervals
$I_1,\ldots,I_M\subset(0,1)$. Let
$\psi_j\in C_c^\infty(I_j)$ satisfy $0\leq\psi_j\leq1$ and
$\max_x\psi_j(x)=1$. Let $\Delta=1/4$.

Choose a smooth cutoff $\chi:[0,1]\to[0,1]$ that equals one at zero and equals
zero on every good and bad state coordinate. For $j\in[M]$, define the smooth
ambient functions
\begin{equation}
 F_{j,+}(a,y)=\frac12+\Delta\chi(y)\psi_j(a),
 \qquad
 F_{j,-}(a,y)=\frac12-\Delta\chi(y)\psi_j(a).
 \label{eq:ambient-functions}
\end{equation}
Their restrictions to $\iota(\cZ)$ have finite $\cH$ norms. Put
\begin{equation}
 A_T=\max\left\{
  \max_{j\in[M],\,\sigma\in\{+,-\}}\|F_{j,\sigma}|_{\iota(\cZ)}\|_{\cH},
  \tfrac12\|1\|_{\cH}
 \right\}<\infty
\end{equation}
and choose the integer
\begin{equation}
 N_T\geq A_T/u.
 \label{eq:clone-number}
\end{equation}
The stage-one transitions of instance $j$ are, for every $z=(s,a)$,
\begin{align}
 P_{1,j}(g_n\mid z)&=\frac{F_{j,+}(a,\eta(s))}{N_T},
 &n&\leq N_T,\label{eq:hard-good}\\
 P_{1,j}(b_n\mid z)&=\frac{F_{j,-}(a,\eta(s))}{N_T},
 &n&\leq N_T.\label{eq:hard-bad}
\end{align}
All remaining target states have probability zero. The two displayed families sum
to one and are nonnegative. At stage two, let the transition be independent of
$z$ and uniform on $\{g_n,b_n:n\leq N_T\}$. These final transitions do not affect
return, but they make the coordinate assumption valid for every $h\in[H]$ under
the convention that $P_H$ is specified.

Equations~\eqref{eq:ambient-functions}--\eqref{eq:clone-number} give
\begin{equation}
 \|P_{h,j}(s'\mid\cdot)\|_{\cH}\leq u
 \qquad\text{for every }h\in\{1,2\},\ s'\in\cS.
 \label{eq:coordinate-verified}
\end{equation}
When starting at $s_0$, the total probability of a good terminal state is
\begin{equation}
 \sum_{n\leq N_T}P_{1,j}(g_n\mid(s_0,a))
 =\frac12+\Delta\psi_j(a).
\end{equation}
The episode's expected return is therefore
$\alpha(1/2+\Delta\psi_j(a))$, whose optimum is
$\alpha(1/2+\Delta)$.

The flat instance $0$ replaces every $\psi_j$ by zero and uses the same $N_T$.
Conditioned on ``good'' or ``bad'', the observed index $n$ is uniform on
$[N_T]$ under every instance. Thus the label contains no information beyond the
Bernoulli outcome.

\subsection{Coupling and regret}

Fix an arbitrary randomized history-dependent algorithm $\mathsf A$. Under the flat
instance, let $A_t$ be its first-stage action in episode $t$. For each $j$, define
\begin{equation}
 \tau_j=\inf\{t\leq T:A_t\in\supp(\psi_j)\},
\end{equation}
with $\inf\varnothing=\infty$. For every realized trajectory of actions,
\begin{equation}
 \sum_{j=1}^{M}\mathbf 1\{\tau_j\leq T\}\leq T
 \label{eq:path-hit-count}
\end{equation}
because the supports are disjoint. Taking expectations under the flat instance
and dividing by $M=2T$ gives
\begin{equation}
 \frac1M\sum_{j=1}^M\Pr_0(\tau_j\leq T)\leq\frac12.
 \label{eq:average-hit}
\end{equation}

For a fixed $j$, couple instance $j$ and the flat instance with the same algorithm
randomness and the same random bits used to draw terminal outcomes. Before an
action enters $\supp(\psi_j)$, the good probability is $1/2$ in both instances,
and the conditional clone label is uniform in both. Their histories and next
actions consequently agree until entry. Entry is determined by the action before
the new outcome is observed. Therefore
\begin{equation}
 \Pr_j(\tau_j\leq T)=\Pr_0(\tau_j\leq T).
 \label{eq:hitting-coupling}
\end{equation}
Equation~\eqref{eq:average-hit} supplies some $j^*$ with
$\Pr_{j^*}(\tau_{j^*}>T)\geq1/2$. On this event every selected action has expected
gap $\alpha\Delta$. The expected realized pseudo-regret equals the standard
value regret by the tower property, so
\begin{equation}
 \E_{j^*}R_T
 =\E_{j^*}\sum_{t=1}^T\alpha\Delta(1-\psi_{j^*}(A_t))
 \geq\frac{\alpha\Delta}{2}T.
 \label{eq:lower-complete}
\end{equation}
The argument already includes the algorithm's private randomness and proves the
minimax claim with $c_\nu=\alpha\Delta/2$.

Finally, every Gram matrix of $k$ is a Gram matrix of the ambient Mat\'ern kernel on
$[0,1]^2$. Hence
\begin{equation}
 \gamma_T(k;1)\leq\gamma_T(\kappa_\nu;1)
 =\widetilde O\!\left(T^{2/(2\nu+2)}\right)=o(T),
\end{equation}
using the standard Mat\'ern information-gain bound
\citep{vakili2021information}. This completes the proof of
Theorem~\ref{thm:linear}.

\section{Proofs and separations for Bellman radii}
\label{app:radius}

\subsection{Proof of Theorem~\ref{thm:radius}}

For any state $y$, take $V=\mathbf1_{\{y\}}$ in
equation~\eqref{eq:radii}. This gives $\Crkhs\leq\Bpos$. The inclusion of the
positive unit cube in the signed unit ball gives $\Bpos\leq\Bsgn$.
For any $V$ with $\|V\|_\infty\leq1$, write $V=V_+-V_-$, where
$0\leq V_+,V_-\leq1$. Then
\begin{equation}
 \|PV\|_{\cH}\leq\|PV_+\|_{\cH}+\|PV_-\|_{\cH}
 \leq2\Bpos(P;\cH).
\end{equation}
For $0\leq V\leq1$ on a countable space, the triangle inequality gives
\begin{equation}
 \|PV\|_{\cH}
 =\left\|\sum_yV(y)p_y\right\|_{\cH}
 \leq\sum_yV(y)\|p_y\|_{\cH}
 \leq\Vrkhs(P;\cH).
\end{equation}
Taking suprema proves equation~\eqref{eq:hierarchy}.

For cloning, let
$\bar V(y)=m_y^{-1}\sum_{i=1}^{m_y}\widetilde V(y,i)$. Direct substitution gives
\begin{equation}
 \widetilde P\widetilde V=P\bar V.
 \label{eq:average-value}
\end{equation}
The average of a positive or signed unit-ball function is in the corresponding
unit ball. Conversely, every original function has a constant-on-fibers lift.
The suprema defining both radii are equal. Also
\begin{equation}
 \Vrkhs(\widetilde P;\widetilde\cH)
 =\sum_y\sum_{i=1}^{m_y}\frac{\|p_y\|_{\cH}}{m_y}
 =\Vrkhs(P;\cH).
\end{equation}

For an $\cH$-valued countably additive measure
$\mathsf m(A)=P(A\mid\cdot)$, its semivariation on $\cS$ can be written
\begin{equation}
 \|\mathsf m\|_{\mathrm{sv}}(\cS)
 =\sup_{\|V\|_\infty\leq1}
 \left\|\int_{\cS}V\,d\mathsf m\right\|_{\cH}.
\end{equation}
This is exactly $\Bsgn$. On countable spaces, countable additivity in $\cH$ follows
whenever the slice variation is finite. More generally, the identity applies
whenever the transition set function is an $\cH$-valued vector measure. The
definition of $\Bpos$ directly proves its minimal-constant claim.

\subsection{No dimension-free reverse inequalities}

First let the current domain contain one point and let its RKHS be $\R$ with the
usual norm. A transition uniform on $n$ next states has
\begin{equation}
 \Crkhs=1/n,
 \qquad \Bpos=\Bsgn=\Vrkhs=1.
\end{equation}
Thus $\Bpos/\Crkhs=n$.

Second let the current domain be $[n]$ with the Kronecker kernel
$k(i,j)=\mathbf1\{i=j\}$. Let the transition be deterministic, sending current
point $i$ to next state $i$. Its coordinate vectors form the Euclidean basis, so
\begin{equation}
 \Crkhs=1,
 \qquad \Bpos=\Bsgn=\sqrt n,
 \qquad \Vrkhs=n.
\end{equation}
Hence both remaining gaps can diverge.

\subsection{Proof of Proposition~\ref{prop:density}}

For $0\leq V\leq1$, Bochner integrability and the triangle inequality imply
\begin{align}
 \|PV\|_{\cH}
 &=\left\|\int_{\cS}V(s')p(s'\mid\cdot)\,d\mu(s')\right\|_{\cH}\\
 &\leq\int_{\cS}V(s')\|p(s'\mid\cdot)\|_{\cH}\,d\mu(s')\\
 &\leq\Vrkhs_\mu(P;\cH).
\end{align}
If the pointwise norm is at most $u$, the last expression is at most
$u\mu(\cS)$.

Suppose $d\mu'=w\,d\mu$ with $w>0$ almost everywhere. The transition density
relative to $\mu'$ is $p'=p/w$, so
\begin{equation}
 \int\|p'(s'\mid\cdot)\|_{\cH}\,d\mu'(s')
 =\int\frac{\|p(s'\mid\cdot)\|_{\cH}}{w(s')}w(s')\,d\mu(s')
 =\Vrkhs_\mu(P;\cH).
\end{equation}
A measurable bijective relabeling gives the same identity by change of variables.

\subsection{Proof of Corollary~\ref{cor:closure}}

For $h=H$, the claim is immediate because $V=0$. If $h<H$ and
$0\leq V\leq H-h$, then $V/(H-h)$ belongs to the positive unit ball. Therefore
\begin{equation}
 \|P_hV\|_{\cH}\leq(H-h)\Bpos(P_h;\cH).
\end{equation}
The triangle inequality after adding $r_h$ proves
equation~\eqref{eq:closure}. Setting $r_h=0$ and taking a supremum over $V$ proves
minimality.

\section{Detailed proof-level literature audit}
\label{app:audit}

This appendix records the exact distinctions summarized in
Table~\ref{tab:audit}. It is not a claim that every theorem in a cited paper fails.
The issue is one implication used to instantiate a Bellman-image premise.

\paragraph{Direct Bellman closure.}
\citet[Assumption 4.1]{yang2020kernel} assumes that the Bellman operator maps every
bounded action-value function into a fixed RKHS ball. Their main kernel theorem is
stated under that assumption and is not affected. Immediately afterward, their
equation (4.1) gives coordinate transition slices in the unit ball as a sufficient
condition and asserts an $H+1$ target norm. Theorem~\ref{thm:radius} shows that this
sufficiency statement needs a bound on $\Bpos$, slice variation, or state volume.

\paragraph{The source of the volume term.}
\citet[Assumption 2]{yeh2023sample} assumes a pointwise RKHS norm bound for the
transition density. Their Lemma 3 proves
\begin{equation}
 \|PV\|_{\cH}\leq\frac{c}{1-\gamma},
 \qquad c=\int_{\cS}ds',
\end{equation}
for discounted values bounded by $1/(1-\gamma)$. Their appendix proof displays the
integral and the factor $c$. This agrees with Proposition~\ref{prop:density}. The
lemma is valid when the density convention and finite $c$ are part of the model.

This distinction also identifies a valid normalized special case.
\citet[Lemma 2.3]{vakili2024rewardfree} states its value-function domain as
$[0,1]^d$ and invokes the same density calculation. Under the stated Lebesgue
unit-cube convention, $c=1$, so the displayed radius-$H$ conclusion retains the
entire volume factor. An extension to other regular compact domains, as suggested
later in that paper, must retain their volume or an invariant integrated bound.

\paragraph{Dimension-free uses.}
\citet[Assumption 1 and Lemma 1]{vakili2023kernel} assumes unit coordinate norms and
calls membership of every $r_h+P_hV$ in the radius-$(H+1)$ ball an immediate
consequence, citing the preceding lemma. The cited proof contains $c$. Thus the
dimension-free consequence is not established on the paper's general state space.
Theorem~\ref{thm:linear} applies directly to the transition-only Assumption 1 posed
by \citet{vakili2024open}.

\citet[Assumptions 2 and 4]{vakili2024average} separately assumes coordinate
transition norms and optimistic closure for value functions in a state RKHS. In
the proof of its Theorem 3, the required state-action RKHS bound
$\|Pv\|_{\cH}=O(w)$ is then obtained from the coordinate premise by citing Yeh et
al.'s Lemma 3. Because its state space is general, that step requires the omitted
reference-volume or integrated-variation constant. The paper's self-contained
confidence theorem is stated directly with a bound on $\|Pv\|_{\cH}$ and is not
affected by this point.

\citet[Assumption 1]{kayal2025rewardfree} likewise states
$\|PV\|_{\cH}=O(H)$ and cites Yeh et al. The statement is correct only if the
suppressed constant includes the reference volume or an integrated variation
bound. The same distinction applies to the coordinate premise and imported KOVI
bound in \citet{cioba2026invariances}, and to the arbitrary-policy $Q$-norm claim
following equation (4) of \citet{pavlovic2026preferences}. These are recent
preprints or applications with additional assumptions. We make no claim here about
parts of their arguments that do not use the coordinate-to-Bellman step.

\paragraph{Independent issues and invariant alternatives.}
Appendix F of \citet{maran2024local} observes that a cover restricted to a current
partition cell does not control values at unrestricted next states in
\citet{vakili2023kernel}. That moving-target localization issue is independent of
the coordinate-volume obstruction proved here. The smooth-Bellman assumptions and
local linearity results of \citet{maran2024local,maran2024smooth} directly control
Bellman images and survive cloning.

\citet{chowdhury2023cme} assumes an operator-level conditional mean embedding that
maps a chosen state RKHS into the state-action RKHS, together with closure of the
algorithm's optimistic value class. Such operator assumptions are invariant under
consistent pullbacks. They are not identical to $\Bpos$, since their domain is a
specified state RKHS rather than all bounded functions. The measure-theoretic CME
literature carefully distinguishes these representability conditions
\citep{park2020measure,hertel2026regularity}.

The 2025 OpenReview preprint of \citet{kumar2025closure} also assumes a bounded
conditional mean embedding and projects optimistic proxies into a chosen state-RKHS
ball. This is a recent algorithmic use of the invariant operator route. It neither
derives Bellman closure from coordinate transition slices nor studies state cloning,
so its contribution is complementary to ours.

\section{Verification artifact}
\label{app:checks}

The file \texttt{audit\_cloning.py} performs four independent finite checks.
\begin{enumerate}[leftmargin=*,itemsep=2pt,topsep=2pt]
 \item It samples a transition matrix on a finite current domain, computes exact
 finite-RKHS coordinate norms, and verifies their prescribed clone scaling.
 \item It enumerates the vertices of the positive and signed value cubes to compute
 both Bellman radii. It verifies that they are unchanged after nonuniform cloning.
 \item It verifies invariance of slice variation and equality of original and
 pulled-back information gains for many point sequences.
 \item It exhaustively checks the support-count inequality
 \eqref{eq:path-hit-count} for every length-four sequence over eight supports and
 verifies the resulting one-half no-hit fraction.
\end{enumerate}
All comparisons use a tolerance of $10^{-10}$. These checks are not substitutes for
the proofs. They guard against indexing, scaling, and inequality-direction errors.
